\section{Introduction}
\label{sec:introduction}

A medical answer depends on the evidence that supports it. When a patient's condition changes, a system may need to add a newly indicated treatment, withdraw an option whose prerequisites no longer hold, or reconsider the leading diagnosis. Other changes should leave the answer intact. Selective updating requires tracking how the changed information bears on the particular judgment being requested.

Consider a treatment supported by two independent indications. Removing one indication leaves the other available to justify treatment; withdrawing the treatment requires checking both support paths. A diagnostic update requires a different operation: new evidence can strengthen a candidate while another diagnosis remains more likely. The final answer's change or preservation captures the outcome of these decisions. The underlying operation---adding support, removing support, or comparing alternatives---explains what the system must compute.

Retrieval-augmented generation (RAG) gives medical language models access to external knowledge, and agent-based systems add iterative retrieval, multiple answers, or explicit control of reasoning \citep{xiong2024medrag,li2026medrgag,wu2026marag,jiang2025tcrag}. Using that evidence requires an appropriate decision procedure for the requested update. Existing counterfactual, perturbation, and robustness resources test complementary aspects of this ability through different units, answer spaces, and evaluation rules. A shared behavioral vocabulary can connect their observations while preserving their original tasks.

We introduce a taxonomy of five functional relations between a change and a target judgment: support addition, support removal, candidate comparison, consequence calculation, and target preservation. The categories can overlap: the same change may add evidence for a candidate and require a new comparison among alternatives. We use this taxonomy to organize an inventory of 26 candidate resources and construct a benchmark spanning 11 sources. Its 6,104 evaluation units retain the native questions, options, target definitions, and source-specific scoring definitions.

Functional evaluation reveals weaknesses hidden by aggregate rankings. Seven configurations from six system families achieve aggregate scores between 47.35 and 51.61, while their scores on support removal range from 7.59 to 41.77. Different systems lead on different functional categories. These profiles motivate an operation-centered approach: identify the computation required by the observed failure mode, then attach a component that performs it through the existing interface.

We organize these components as a counterfactual mixture of agents (\cfmoa{}), with five functional roles. Its support agents execute three-valued rules, its comparison agent makes an explicit decision over a diagnosis catalog, its consequence agent applies local interventions to a supplied executable model, and its preservation agent stabilizes native answers through structured readout. A sparse controller selects the applicable operation from the visible input and available capabilities. This design gives each agent a concrete responsibility and allows compatible operations to share an implementation.

Our contributions are threefold. First, we introduce a functional taxonomy and a provenance-preserving benchmark that connect diverse medical counterfactual and robustness tasks. Second, we provide a seven-configuration diagnostic evaluation, together with head-level comparisons that improve all 21 system--category combinations on support addition, support removal, and candidate comparison. Third, we develop \cfmoa{} and evaluate the complete adopted controller on two MedRAG configurations, obtaining aggregate gains of 10.09 and 12.01 percentage points and improved correctness on both changing and preserving pairs. The accompanying controls explain which operations account for these improvements and when further aggregation adds little value.
